% Swarm-blocker 6, v2 / Agent 07 of N
% Voices: Kontsevich + Caldararu + Markarian + Costello-Li
% Task: HKR comparison of HH^*(D^b(K3 x E)), PV^*(K3 x E),
%   and the hCS gauge dg-Lie Omega^{0,*}(X, g)[1].
% Primary anchors:
%   Hochschild-Kostant-Rosenberg 1962 Trans.\ AMS 102 (HKR original);
%   Kontsevich 1999 Lett.\ Math.\ Phys.\ 48 (deformation quantisation,
%     Tsygan-Kontsevich formality);
%   Markarian 2009 J.\ Lond.\ Math.\ Soc.\ 79 (HKR for D^b(X) with
%     Mukai-pairing twist);
%   Caldararu 2003 Adv.\ Math.\ 194 (Mukai pairing on D^b);
%   Caldararu-Willerton 2010 New York J.\ Math.\ 16 (Mukai pairing,
%     Hodge identification);
%   Bondal-Orlov 2002 Compositio Math.\ 125 (FM kernels, autoequivalences);
%   Calaque-Van den Bergh 2010 Adv.\ Math.\ 224 (cochain-level HKR
%     globalisation, Duflo square-root td(X)^{1/2});
%   Costello-Li 2012, 2016 (BV propagator for hCS, BCOV one-loop
%     curving alpha_BCOV);
%   Kapranov 1999 Compositio Math.\ 115 (L_infty on T_X[-1]).

\providecommand{\HH}{\mathrm{HH}}
\providecommand{\PV}{\mathrm{PV}}
\providecommand{\hCS}{\mathrm{hCS}}
\providecommand{\Obs}{\mathrm{Obs}}
\providecommand{\Coh}{\mathrm{Coh}}
\providecommand{\Perf}{\mathrm{Perf}}
\providecommand{\BCOV}{\mathrm{BCOV}}
\providecommand{\HKR}{\mathrm{HKR}}
\providecommand{\PhiFA}{\Phi^{\mathrm{FA}}}
\providecommand{\Bun}{\mathrm{Bun}}
\providecommand{\MC}{\mathrm{MC}}
\providecommand{\dbar}{\bar\partial}
\providecommand{\frakg}{\mathfrak{g}}
\providecommand{\Ethree}{E_3}
\providecommand{\Eone}{E_1}
\providecommand{\BV}{\mathrm{BV}}
\providecommand{\C}{\mathbb{C}}
\providecommand{\Z}{\mathbb{Z}}
\providecommand{\sym}{\mathrm{Sym}}
\providecommand{\End}{\mathrm{End}}
\providecommand{\cO}{\mathcal{O}}
\providecommand{\cC}{\mathcal{C}}
\providecommand{\ClaimStatusConditional}[1][]{\textbf{[conditional#1]}}

\section*{HKR comparison of $\HH^\bullet(D^b(K3\times E))$,
 $\PV^\bullet(K3\times E)$ and the hCS gauge dg-Lie}
\label{sec:swarm-blocker6-v2-agent07-hh-pv-comparison}

\paragraph{The blocker.}
The Stage-$1$ factorisation construction $\PhiFA_3(D^b\Coh(X))$ feeds
on the $\Ethree$-algebra $\HH^\bullet(D^b\Coh(X))$. The geometric input
to holomorphic Chern--Simons on $X = K3\times E$ is the gauge dg-Lie
algebra
\[
 L_\frakg^X \;:=\; \Omega^{0,\bullet}(X, \frakg)[1],
 \qquad d = \dbar, \qquad
 [\alpha\otimes a,\;\beta\otimes b] = \alpha\wedge\beta\,\otimes\,[a,b]_\frakg.
\]
At the level of cohomology, $H^\bullet(L_\frakg^X) = \frakg\otimes
H^{0,\bullet}(X)$ has no $T_X$-factor; HKR gives
$\HH^\bullet(D^b\Coh(X)) \simeq \PV^\bullet(X) = \bigoplus_{p,q}
H^q(X,\wedge^p T_X)$, which carries $T_X$ explicitly. The two grammars
are not in one-to-one correspondence pointwise. The blocker is to name
the precise shape of the comparison so that $\PhiFA_3(D^b\Coh(X))$ does
not over-count or under-count the geometric content.

The resolution is a boundary, not an equality. $L_\frakg^X$ is the
gauge dg-Lie at a \emph{fixed point} of the holomorphic bundle moduli
problem. The HKR target $\PV^\bullet(X)$ is the categorical
Hochschild/polyvector complex of $\Perf(X)$. Dolbeault forms with
coefficients in a finite Lie algebra resolve holomorphic
$\frakg$-valued functions; they do not resolve all polyvectors. The
polyvector grammar meets field theory through BCOV/Kodaira--Spencer
theory, whose field complex uses $T_X[-1]$, and through a separate
framed hCS--Hochschild primitive package. Substituting
$\frakg=T_X[-1]$ changes the theory; it is not ordinary hCS with a
different finite gauge algebra.

\subsection*{(a) HKR for $D^b(X)$, $X$ smooth proper}

\begin{theorem}[Hochschild--Kostant--Rosenberg for $D^b\Coh(X)$]
\label{thm:hkr-Db-X}
Let $X$ be a smooth proper variety over $\C$ of dimension $d$. The
HKR map
\[
 I^{\HKR}_X \colon\;
 \HH^\bullet(D^b\Coh(X))
 \;\xrightarrow{\sim}\;
 \PV^\bullet(X)
 \;:=\; \bigoplus_{p+q=\bullet}
 H^q\bigl(X,\,\wedge^p T_X\bigr)
\]
is a quasi-isomorphism of complexes
(\emph{Hochschild--Kostant--Rosenberg} 1962;
\emph{Markarian} 2009 J.\ Lond.\ Math.\ Soc.\ 79
\textup{[}arXiv:math/$0610553$\textup{]} for the $D^b$-version with
Mukai twist;
\emph{C\u ald\u araru} 2003 Adv.\ Math.\ 194
for the Mukai pairing). The Caldararu--Willerton \emph{Mukai-twisted
HKR} map
\[
 I^{\HKR,\mathrm{CW}}_X \;:=\; I^{\HKR}_X \circ \mathrm{td}(X)^{1/2}
\]
is the unique quasi-isomorphism upgrading the underlying
chain-quasi-isomorphism to a Gerstenhaber morphism on cohomology
(\emph{Calaque--Van den Bergh} 2010 Adv.\ Math.\ 224); on a Calabi--Yau
$X$ ($c_1(X)=0$) the Todd correction reduces to the Duflo cocycle
$\mathrm{td}(X)^{1/2} = 1 + c_2(T_X)/24$ (no $c_1$-piece).
\end{theorem}

The completion appropriate to a smooth proper $X$: $\PV^\bullet(X)$ is
\emph{finite-dimensional in each cohomological degree} (each
$H^q(X,\wedge^p T_X)$ is finite-dimensional by Serre's theorem on
coherent cohomology of a proper variety); the HKR map respects this
finiteness without further completion. For non-proper $X$ (open
varieties, formal disks, $\C^d$) the polyvector ring is replaced by
the Dolbeault polyvector complex
$\Omega^{0,\bullet}(X, \wedge^\bullet T_X)$ with $\dbar$-differential,
and the corresponding HKR statement requires a Hochschild-cochain
completion (see Calaque--Van den Bergh op.\ cit.\ \S$3$).

\subsection*{(b) HKR-K\"unneth for $X = K3\times E$}

\begin{proposition}[K\"unneth decomposition of $\PV^\bullet(K3\times E)$]
\label{prop:pv-kunneth-k3xe}
For $X = K3\times E$, the polyvector cohomology factors:
\[
 H^q(X, \wedge^p T_X)
 \;=\;
 \bigoplus_{p_1+p_2 = p,\ q_1+q_2 = q}
 H^{q_1}(K3, \wedge^{p_1} T_{K3})
 \;\otimes\;
 H^{q_2}(E, \wedge^{p_2} T_E),
\]
because $T_{K3\times E} = \pi_1^* T_{K3} \oplus \pi_2^* T_E$ and
exterior powers split:
$\wedge^p T_X = \bigoplus_{p_1+p_2=p}
 \pi_1^*\wedge^{p_1}T_{K3}\otimes\pi_2^*\wedge^{p_2}T_E$.
\end{proposition}

\begin{proof}
For $X_1, X_2$ smooth proper, $T_{X_1\times X_2}$ splits canonically
along the projections; $\wedge^p$ commutes with direct sums of vector
bundles via the binomial expansion; K\"unneth in coherent cohomology
holds for smooth proper varieties (Hartshorne III.6).
\end{proof}

\paragraph{$\PV^\bullet(K3)$.}
$T_{K3}$ has rank $2$; $\wedge^2 T_{K3} = \det T_{K3} = K_{K3}^{-1} = \cO_{K3}$
by the Calabi--Yau condition $K_{K3} = \cO_{K3}$; $\wedge^0 T_{K3} = \cO_{K3}$.
The Dolbeault cohomology of $\cO_{K3}$ is $H^{0,\bullet}(K3) = (1, 0, 1)$
in degrees $(0, 1, 2)$. The cohomology of $T_{K3}$ is, by Serre duality
$H^q(K3, T_{K3}) \simeq H^{2-q}(K3, \Omega^1_{K3})^\vee$, equal to
$(0, 20, 0)$ in degrees $(0,1,2)$ since $h^{1,1}(K3) = 20$ and
$h^{0,1}(K3) = h^{0,1}(K3)^\vee$ via Serre $K_{K3} = \cO$. Bidegrees:
\[
\begin{array}{c|ccc}
 \PV^\bullet(K3) & p=0\ (\cO) & p=1\ (T_{K3}) & p=2\ (\wedge^2 T_{K3}=\cO) \\\hline
 q=0 & 1 & 0 & 1 \\
 q=1 & 0 & 20 & 0 \\
 q=2 & 1 & 0 & 1
\end{array}
\]
Total dimension $\dim_\C \PV^\bullet(K3) = 1+0+1+0+20+0+1+0+1 = 24$,
which is the rank of the Mukai lattice $\widetilde H^\bullet(K3,\Z)$.
This is no accident: the Caldararu--Willerton Mukai pairing identifies
$\HH^\bullet(D^b\Coh(K3)) \simeq \widetilde H^\bullet(K3,\C)$ via
twisted HKR (\emph{C\u ald\u araru--Willerton} 2010, \S$5$).

\paragraph{$\PV^\bullet(E)$.}
$E$ is a complex elliptic curve, $T_E = \cO_E$, $\wedge^0 T_E = \cO_E$;
no higher exterior powers. $H^{0,\bullet}(E) = (1, 1)$ in degrees $0, 1$.
\[
\begin{array}{c|cc}
 \PV^\bullet(E) & p=0\ (\cO) & p=1\ (T_E=\cO) \\\hline
 q=0 & 1 & 1 \\
 q=1 & 1 & 1
\end{array}
\]
Total $\dim_\C\PV^\bullet(E) = 4$.

\paragraph{Product.}
By Proposition~\ref{prop:pv-kunneth-k3xe},
$\dim_\C\PV^\bullet(K3\times E) = 24\cdot 4 = 96$. The bidegree
$(p,q)$-decomposition is the K\"unneth tensor product of the two
tables. As a $\Z$-graded vector space (collapsing to $p+q$), one
checks: $\sum_n h_n t^n = (\sum_{p,q} h^{p,q}_{K3}t^{p+q})\cdot
(\sum_{p,q}h^{p,q}_E t^{p+q}) = (1+t)^2(1+22t^2+t^4) /
\textit{(check)}$ -- the explicit polynomial is what matters for the
HKR comparison below.

\subsection*{(c) The hCS gauge dg-Lie cohomology on $X = K3\times E$}

The hCS gauge dg-Lie at the trivial bundle is
\[
 L_\frakg^X \;=\; \Omega^{0,\bullet}(X, \frakg)[1],
 \qquad H^\bullet(L_\frakg^X) \;=\; \frakg\otimes H^{0,\bullet}(X)[1].
\]
For $X = K3\times E$ with the K\"unneth Hodge diamond
\[
 H^{0,\bullet}(X) \;=\;
 H^{0,\bullet}(K3)\otimes H^{0,\bullet}(E)
 \;=\; (1,0,1)\otimes(1,1) \;=\; (1,1,1,1)
\]
in degrees $0,1,2,3$, so
\[
 H^\bullet(L_\frakg^X)
 \;=\;
 \frakg^{(0)}\oplus\frakg^{(1)}\oplus\frakg^{(2)}\oplus\frakg^{(3)},
\]
each summand of dimension $\dim\frakg$, total $4\dim\frakg$.
For $\frakg = \mathfrak{sl}_2$: $4\cdot 3 = 12$. (Cf.\
\texttt{cy\_to\_chiral.tex} Eq.~\eqref{eq:harmonic-basis-k3xe}
display, which fixes the harmonic basis
$1, d\bar z_E, \bar\omega_{K3}, \bar\omega_{K3}\wedge d\bar z_E$.)

\subsection*{(d) Apparent mismatch}

\begin{itemize}
\item $\HH^\bullet(D^b\Coh(X)) \simeq \PV^\bullet(X) =
 \bigoplus_{p,q} H^q(X,\wedge^p T_X)$ \emph{has $T_X$-factors};
 dimension $96$ on $X = K3\times E$.
\item $H^\bullet(L_\frakg^X) = \frakg\otimes H^{0,\bullet}(X)$
 \emph{has only $\cO_X$-factors}; dimension $4\dim\frakg$.
\end{itemize}
The two are not isomorphic, even after tensoring with $\frakg$. The
hCS cohomology has \emph{no $T_X$-factor}, while the HKR polyvector
target carries all of $\wedge^\bullet T_X$.

This is not a flaw: the hCS gauge dg-Lie is the perturbative observable
algebra at \emph{a single bundle} (the trivial one), while
$\HH^\bullet(D^b\Coh(X))$ is a categorical invariant of the
\emph{whole} category $D^b\Coh(X)$, which contains every coherent
sheaf and every Fourier--Mukai autoequivalence. The $T_X$-factor
records geometric polyvector directions of the category; it is
invisible at a fixed-bundle perturbative expansion.

\subsection*{(e) Resolution: a typed comparison boundary}

\begin{theorem}[Typed HKR--hCS--BCOV comparison boundary]
\label{thm:hkr-vs-hcs-resolution}
\ClaimStatusConditional{\textup{(}comparison boundary; no
quasi-isomorphism between ordinary hCS and $\PV^\bullet$\textup{)}}
Let $X$ be a smooth proper Calabi--Yau threefold and $G$ a reductive
group. Let $\Bun_G(X)$ denote the derived moduli stack of holomorphic
$G$-principal bundles on $X$ (a derived Artin stack; \emph{Toen--Vaqui\'e}
2007 Ann.\ Sci.\ ENS 40). Then:
\begin{enumerate}[label=(\roman*)]
\item the ordinary hCS perturbative complex at a fixed bundle has
linear fields $L_\frakg^X=\Omega^{0,\bullet}(X,\frakg)[1]$ and
cohomology $\frakg\otimes H^{0,\bullet}(X)[1]$;
\item the categorical Hochschild cohomology of $\Perf(X)$ is identified
by HKR with $\bigoplus_{p,q}H^q(X,\wedge^pT_X)$, with the
Caldararu--Willerton/Todd correction governing multiplicative and
Gerstenhaber compatibility;
\item the BCOV/Kodaira--Spencer theory has fields built from
$\Omega^{0,\bullet}(X,\PV^\bullet_X)$ \textup{(}at first tangent layer,
$\Omega^{0,\bullet}(X,T_X[-1])$\textup{)} and is the field theory
naturally tied to polyvectors;
\item an hCS-to-Hochschild comparison is therefore a span
\[
 \Obs^{q,\hCS}_{P_0}(X,\frakg)
 \xleftarrow{\;\mathrm{framed\ primitive}\;}
 \mathcal N_{X,\frakg}^{\mathrm{HH/BV}}
 \xrightarrow{\;\mathrm{HKR/BCOV}\;}
 \PhiFA_3(\Perf(X)) ,
\]
not a quasi-isomorphism
$\Omega^{0,\bullet}(X,\frakg)\simeq \frakg\otimes\PV^\bullet(X)$.
\end{enumerate}
\end{theorem}

\begin{proof}
(i) is the standard BV-perturbative content: at a stationary point of
the BV action, the linearised observable complex is built from the
gauge dg-Lie shifted by $1$
(Costello--Gwilliam vol.\ $1$, Ch.\ $5$); for hCS at the trivial
bundle, the gauge dg-Lie is $L_\frakg^X = \Omega^{0,\bullet}(X,\frakg)[1]$.

(ii) is HKR for $\Perf(X)$ with the standard Mukai/Todd correction. It
concerns the diagonal of $X$ and the category of perfect complexes, not
the Dolbeault resolution of a fixed principal $G$-bundle.

(iii) is Kapranov 1999 Compositio Math.\ 115 \S$4$ together with the
BCOV field content: the polyvector sheaf carries the
Schouten--Nijenhuis bracket and the Atiyah-sourced higher
$L_\infty$ operations. This is a different BV theory from ordinary hCS
with finite gauge Lie algebra $\frakg$.

(iv) packages the two theories in a span. The middle object is exactly
the extra framed primitive data: a choice of open--closed source map,
zero-mode splitting, completion, and compatibility homotopies.
\end{proof}

\begin{remark}[Two grammars, two scopes]
\label{rem:two-grammars-scopes}
The two grammars compute distinct things at a fixed bundle:
\begin{itemize}
\item the hCS gauge cohomology $H^\bullet(L_\frakg^X) = \frakg\otimes
 H^{0,\bullet}(X)$ counts \emph{$\frakg$-valued harmonic forms},
 i.e.\ the perturbative spectrum at the trivial bundle;
\item the HKR polyvector $\PV^\bullet(X) = \bigoplus_{p,q}
 H^q(X,\wedge^p T_X)$ counts \emph{infinitesimal sheaf moduli},
 weighted by the $\wedge^p$-exterior power.
\end{itemize}
$\PhiFA_3(D^b\Coh(X))$ uses the second; the hCS perturbative
calculation realises the first as the BV-quantised
\emph{trivial-bundle stratum} of the gauge direction. The Maurer--Cartan
equation
$\dbar A + \tfrac12[A,A] = 0$ on $A\in L_\frakg^X$ in degree zero
parametrises the formal moduli of $\frakg$-bundles near trivial; at
the BCOV/Kodaira--Spencer theory, the first tangent field is
$T_X[-1]$ and the Maurer--Cartan equation governs complex-structure
deformations of $X$. This is the $\PV^{p=1,q=1}$ tangent piece. The
higher $\wedge^pT_X$ pieces belong to the BCOV/polyvector grammar, not
to ordinary hCS with finite gauge algebra $\frakg$.
\end{remark}

\subsection*{Cycle log (3 attack-heal rounds)}

\subsubsection*{Cycle 1 -- Attack (Caldararu).}
The HKR isomorphism on $D^b\Coh(X)$ is a quasi-isomorphism of
\emph{complexes}, but the Gerstenhaber-bracket structure on the LHS
does not match the Schouten--Nijenhuis bracket on the RHS unless
twisted by $\mathrm{td}(X)^{1/2}$. On $X = K3\times E$ with $c_1(X) = 0$,
$\mathrm{td}(X)^{1/2} = 1 + c_2(T_X)/24 + \cdots$ has a non-trivial
$c_2$-piece (since $c_2(K3) = 24$, $c_2(K3\times E) = c_2(K3) +
c_2(E) + 2c_1(K3)c_1(E) = 24$, all from the K3 factor); does this
modify the dimension count?

\paragraph{Heal.}
The Duflo $c_2$-correction modifies only the \emph{multiplicative}
structure of the Mukai pairing on $\HH^\bullet$, not the underlying
graded vector space. The dimension count of $\PV^\bullet(K3\times E)$
as $96 = 24\cdot 4$ is unchanged; the HKR comparison of
\emph{vector spaces} is a strict equality. Bracket compatibility
requires $\mathrm{td}^{1/2}$-twist (Calaque--Van den Bergh 2010
Thm.~$3.5.1$); but the comparison with $\PV^\bullet$ at the level of
graded $\C$-vector spaces is the input to dimension counting in (b).

\subsubsection*{Cycle 2 -- Attack (Markarian).}
The Markarian formulation of HKR for $D^b(X)$ is in terms of
$\HH^\bullet(\Perf(X)) = \mathrm{RHom}_{X\times X}(\Delta_*\cO_X,
\Delta_*\cO_X)$, the derived endomorphisms of the diagonal sheaf.
The HKR map is the Markarian \emph{characteristic class} of the
diagonal: $\HH^\bullet \to \PV^\bullet$ via the Atiyah class of the
universal sheaf. But on $X = K3\times E$, the diagonal sheaf
$\Delta_*\cO_X$ does not split as a tensor product of $K3$-diagonal
and $E$-diagonal sheaves under K\"unneth: $\Delta_*\cO_{X_1\times X_2}
\neq \Delta_{X_1,*}\cO_{X_1}\boxtimes\Delta_{X_2,*}\cO_{X_2}$ as
sheaves on $(X_1\times X_2)\times(X_1\times X_2)$. Does HKR-K\"unneth
fail?

\paragraph{Heal.}
The diagonal sheaf $\Delta_*\cO_{X_1\times X_2}$ \emph{does} split
under the canonical isomorphism
$(X_1\times X_2)\times(X_1\times X_2) \simeq (X_1\times X_1)\times
(X_2\times X_2)$ followed by external tensor:
$\Delta_*\cO_{X_1\times X_2} \simeq \Delta_{X_1,*}\cO_{X_1}\boxtimes
\Delta_{X_2,*}\cO_{X_2}$, where $\boxtimes$ is the external tensor
under the swap of factors. K\"unneth for $\mathrm{RHom}_{X\times X}$
follows: $\HH^\bullet(\Perf(X_1\times X_2)) =
\HH^\bullet(\Perf(X_1))\otimes\HH^\bullet(\Perf(X_2))$, and the HKR
map factors as $I^{\HKR}_{X_1}\otimes I^{\HKR}_{X_2}$. On
$X = K3\times E$ this gives the K\"unneth decomposition of
Proposition~\ref{prop:pv-kunneth-k3xe} at the level of derived
endomorphisms.

\subsubsection*{Cycle 3 -- Attack (Costello--Li).}
The hCS gauge dg-Lie at the trivial bundle has cohomology
$\frakg\otimes H^{0,\bullet}(X)$, dimension $4\dim\frakg = 12$ for
$\mathfrak{sl}_2$. The polyvector target has dimension $96$ for
$X = K3\times E$, independent of $\frakg$. How does
$\PhiFA_3(D^b\Coh(K3\times E))$ \emph{not} have a $\frakg$-dependence?
If the chiral output depends on $\frakg$, where in the HKR target does
$\frakg$ enter?

\paragraph{Heal.}
$\PhiFA_3$ acts on the Calabi--Yau category and its Hochschild
cochains. A gauge Lie algebra $\frakg$ enters only after one adds a
separate framed gauge/matrix target or compares to a field theory with
that gauge algebra. Thus the categorical input is
$\HH^\bullet(\cC)$, identified by HKR with $\PV^\bullet(X)$; the hCS
input is the fixed-bundle complex
$\Omega^{0,\bullet}(X)\otimes\frakg[1]$. They are related by the
comparison span of Theorem~\ref{thm:hkr-vs-hcs-resolution}(iv), not by
an identity on cohomology. If $\frakg$ is replaced by $T_X[-1]$, the
resulting field theory is BCOV/Kodaira--Spencer theory; this is a
change of theory and cannot be used to prove an ordinary hCS
quasi-isomorphism.

\subsection*{Conclusion}

\begin{theorem}[Triple HKR comparison boundary on $X = K3\times E$]
\label{thm:triple-hkr-K3xE}
The three complexes
\[
 \HH^\bullet(D^b\Coh(K3\times E)),
 \quad
 \PV^\bullet(K3\times E),
 \quad
 H^\bullet\bigl(\Omega^{0,\bullet}(K3\times E)\otimes\frakg[1]\bigr)
\]
are pairwise related by:
\begin{enumerate}
\item $\HH^\bullet(D^b\Coh(K3\times E)) \stackrel{\HKR,\mathrm{CW}}{\simeq}
 \PV^\bullet(K3\times E)$, via the Caldararu--Willerton Mukai-twisted
 HKR map; both are $96$-dimensional, with K\"unneth decomposition
 $24\cdot 4 = 96$.
\item The hCS gauge cohomology $\frakg\otimes H^{0,\bullet}(K3\times E)$
 is $4\dim\frakg$-dimensional and has no $T_X$-factor. Passing to
 $T_X[-1]$ is not a substitution inside ordinary hCS; it is the
 BCOV/Kodaira--Spencer comparison, whose first tangent layer is
 $H^{0,\bullet}(K3\times E,T_X[-1])$. The full polyvector grammar
 $\PV^\bullet$ belongs to the categorical HKR/BCOV side.
\item At fixed $\frakg$ (a numerical Lie algebra), the relation
 between hCS-BV cohomology and the categorical $\HH^\bullet$ is via
 the framed comparison span: $\Obs^{\hCS}_{\mathrm{triv}}$ is the
 perturbative gauge theory at one bundle basepoint, while
 $\HH^\bullet(D^b\Coh(K3\times E))$ is the categorical Hochschild
 invariant of $\Perf(K3\times E)$. The two do not coincide without
 the framed primitive and completion data.
\end{enumerate}
\end{theorem}

\begin{proof}
(1) is Theorem~\ref{thm:hkr-Db-X} together with
Proposition~\ref{prop:pv-kunneth-k3xe}.
(2) is the separation of ordinary hCS from the BCOV/polyvector theory
(Theorem~\ref{thm:hkr-vs-hcs-resolution}(iii)-(iv)) restricted to
$X = K3\times E$.
(3) follows from cycles 1-3 above and is an instance of
Theorem~\ref{thm:hkr-vs-hcs-resolution}(i)-(ii).
\end{proof}

\paragraph{The remaining obstruction.}
The Caldararu--Willerton twisted HKR is a quasi-isomorphism of
\emph{$L_\infty$-algebras up to $\mathrm{GRT}_1(\Q)$-action}
(\emph{Willwacher} 2014 Invent.\ Math.\ 200), while the
BCOV/polyvector theory carries the Kapranov--Schouten
$L_\infty$-structure. Comparing either of these to ordinary hCS
requires an additional framed primitive, and bracket compatibility is
corrected by a $\mathrm{td}(X)^{1/2}$-cocycle. On a strict
Calabi--Yau (compact, $c_1 = 0$), the Duflo correction is
$1 + c_2(T_X)/24 + \cdots$; on $X = K3\times E$, this is non-trivial
($c_2(K3\times E) = 24$, the K3 second Chern class). The chain-level
$\PhiFA_3(D^b\Coh(K3\times E))$ on the polyvector grammar therefore
carries a Duflo-twisted $\Ethree$-structure that the naive hCS gauge
calculation does not see; the chiral specialisation
$\SpCh_{\Sigma_2,E}\circ\PhiFA_3$ records this twist as the
$\Phi_{10}/\eta^{24}$-associator factor of the $\Delta_5^2/\eta^{24}$
denominator (cf.\ \texttt{hochschild\_calculus.tex}
Remark on $\HH^\bullet(\mathbf H_{\Delta_5})$).

\paragraph{What this resolves.}
The blocker -- ``where does $T_X$ enter the hCS calculation?'' --
resolves into a separation of theories. Ordinary hCS with finite
gauge algebra $\frakg$ never acquires the full $T_X$ tower. The
$T_X$-direction enters the categorical side through HKR and the field
theory side through BCOV/Kodaira--Spencer theory. A comparison with
ordinary hCS is a framed source bridge, not a substitution rule.

\paragraph{Cross-references inside the manuscript.}
\texttt{cy\_to\_chiral.tex}~\eqref{eq:harmonic-basis-k3xe}
fixes the harmonic basis of $H^{0,\bullet}(K3\times E)$;
\texttt{hochschild\_calculus.tex}
Theorem~\ref{thm:kt-formality-d3} fixes the
$\Ethree$-formality on $\HH^\bullet(D^b\Coh(K3\times E))$;
Theorem~\ref{thm:three-atiyah-cocycles} catalogues the three
Atiyah-sourced cocycles $\ell_2, Y_3, \alpha_{\BCOV}$ which live in
$\PV^\bullet(K3\times E)$ and feed Stage~$1$ of the two-stage
factorisation. The BCOV/Kodaira--Spencer comparison is the polyvector
bridge. The hCS-perturbative gauge calculation of
\texttt{cy\_to\_chiral.tex} \S$2071$-$2270$ is the fixed-bundle gauge
side and requires the framed Hochschild primitive before it can be
compared with the categorical $\HH^\bullet$-input of
\texttt{hochschild\_calculus.tex} \S$382$-$465$.
